\section{The SCPTOP benchmark}
\label{appendix:A}

The benchmark against which 3D-SCPTOP is compared is an extension of the
procurement transport model of \autocite{Diaz-Madronero2014APlanning} to the
closed-loop, two-node setting studied here. Three modifications are made to
that model, as stated in Section \ref{sec:results}: fuzzy parameters are
dropped; the transport capacity restriction is expressed through the volume of
the individual RTIs against the truck volume rather than a maximum container
count; and a product may travel in more than one RTI type.

The point of the comparison is to isolate the effect of the loading
representation, so SCPTOP is made \emph{identical to the transport and
inventory block of 3D-SCPTOP in every other respect}: the same two nodes, the
same closed loop of full and folded RTIs, the same lead time, the same coverage
stock, the same payload cap, the same full-truckload floor and the same
objective. It differs in exactly one place --- where 3D-SCPTOP asks whether the
RTIs can be stacked and placed on the trailer floor, SCPTOP asks only whether
their volumes sum to less than the trailer's.

The notation is that of Section \ref{sec:model}: $Y_{gkt}\in\{0,1\}$ indicates
that truck $k$ is dispatched in direction $g$ in period $t$, and
$\overline{V}_{knt}$ and $V^{\mathrm{e}}_{krt}$ are the full RTIs of product
$n$, respectively the folded RTIs of type $r$, that it carries. The symbols
are deliberately those of Section \ref{sec:model}, because the quantity is the
same one and \eqref{eq:A-ftl} below is then literally
\eqref{sb:eq:util-p}--\eqref{sb:eq:utilw-p}; the only difference is that
$\overline{V}_{knt}$ is a decision variable here rather than the slot sum
\eqref{sb:eq:agg-Vk}, SCPTOP having no slots to sum over. The aggregates
$\overline{V}_{nt}$ and $\overline{V^{\mathrm{e}}}_{rt}$ are
\eqref{sb:eq:agg-V}--\eqref{sb:eq:agg-Ve} as before. All remaining notation is
that of Table \ref{tab:sb-Nomenclature}.

\begin{align}
\text{Min } z
&= \sum_{g\in G}\sum_{k\in K}\sum_{t\in T} cK\,Y_{gkt}
 + \sum_{g\in G}\sum_{n\in N}\sum_{t\in T} ic_{gn}I_{gnt}
 \label{eq:A-obj}
\end{align}
subject to, for all $k\in K$ and $t\in T$,
\begin{align}
\sum_{n\in N}\gamma_{p\rho(n)}\overline{V}_{knt} &\le \Gamma_{K}\,Y_{pkt},
 & \sum_{r\in R}\gamma_{sr}V^{\mathrm{e}}_{krt} &\le \Gamma_{K}\,Y_{skt}
 \label{eq:A-vol}\\
\sum_{n\in N}w_{n}\overline{V}_{knt} &\le pK^{w}Y_{pkt},
 & \sum_{r\in R}w^{\mathrm{f}}_{r}V^{\mathrm{e}}_{krt} &\le pK^{w}Y_{skt}
 \label{eq:A-weight}
\end{align}
together with the full-truckload floor on the binding resource, imposed exactly
as in \eqref{eq:col-truck}: a dispatched truck must reach a fraction $\alpha$
of \emph{either} the trailer volume or the payload,
\begin{align}
\sum_{n}\gamma_{p\rho(n)}\overline{V}_{knt} \;\ge\; \alpha\Gamma_{K}\bigl(Y_{pkt}-1+\theta_{pkt}\bigr),
\qquad
\sum_{n}w_{n}\overline{V}_{knt} \;\ge\; \alpha\,pK^{w}\bigl(Y_{pkt}-\theta_{pkt}\bigr),
\label{eq:A-ftl}
\end{align}
with $\theta_{gkt}\in\{0,1\}$ and the supplier direction analogous. The
inventory balances, coverage stock and domains are those of
\eqref{sb:eq:inv-plant-main}--\eqref{sb:eq:safety-empty}, unchanged, and they
read the same two aggregates. The
truck-symmetry row \eqref{dec:sym-truck} is carried here as well, for
the same reason and with the same limited effect: the trucks of a
direction and period are anonymous in both models, and ordering their
\emph{use} leaves the assignment of loads among the used trucks free.
Neither compact model is therefore held to an easier standard than the
other, and the symmetry both retain is what the reformulation of
Section~\ref{sec:decomp} removes by aggregating the blocks.

Constraint \eqref{eq:A-vol} is the whole of the benchmark's loading model: a
truck is deemed able to carry any set of RTIs whose volumes fit inside the
trailer's. No footprint, stack height, layer count or crush limit appears.
Section \ref{sec:results} reports what that omission costs --- not in solution
quality, where SCPTOP appears cheaper, but in whether the plans it prescribes
can be loaded at all.

\paragraph{Why the floor is imposed on the binding resource.}
Writing \eqref{eq:A-ftl} on volume alone, as the original formulation does,
makes dense freight structurally infeasible: filling the fraction
$\alpha=0.55$ of a $97.92$~m$^3$ trailer with freight at $579$~kg/m$^3$ weighs
$31.2$~t against a $24$~t cap, so no truck is dispatchable and the instance is
empty (Section~\ref{sec:model}, where the density is traced to the catalogue). Summing the two ratios
instead is strictly weaker --- it accepts a truck $35\%$ full by volume and
$20\%$ by weight. The disjunctive form above is used identically in every model
compared here, so no method is held to an easier standard than another.
